\documentclass[10pt,a4paper,twocolumn]{article}
\usepackage[utf8]{inputenc}
\usepackage[T1]{fontenc}
\usepackage{amsmath, amssymb}
\usepackage{listings}
\usepackage{xcolor}
\usepackage{geometry}
\geometry{left=1cm,right=1cm,top=1cm,bottom=1.5cm}

\lstset{
    language=C++,
    basicstyle=\ttfamily\scriptsize,
    keywordstyle=\color{blue},
    commentstyle=\color{gray},
    stringstyle=\color{red},
    breaklines=true,
    showstringspaces=false,
    tabsize=2,
    numbers=none,
    extendedchars=true,
    literate={á}{{\'a}}1 {ã}{{\~a}}1 {é}{{\'e}}1 {ê}{{\^e}}1 {í}{{\'i}}1 {ó}{{\'o}}1 {õ}{{\~o}}1 {ô}{{\^o}}1 {ú}{{\'u}}1 {ç}{{\c{c}}}1 {Á}{{\'A}}1 {Ã}{{\~A}}1 {É}{{\'E}}1 {Ê}{{\^E}}1 {Í}{{\'I}}1 {Ó}{{\'O}}1 {Õ}{{\~O}}1 {Ô}{{\^O}}1 {Ú}{{\'U}}1 {Ç}{{\c{C}}}1
}

\begin{document}
\twocolumn
\tableofcontents
\newpage

\section{Introdução}
\begin{lstlisting}
#include <bits/stdc++.h>
using namespace std;

#define fastio ios_base::sync_with_stdio(0); cin.tie(0)
#define all(a) a.begin(), a.end()
#define rall(a) a.rbegin(), a.rend()
#define sz(a) (int) (a).size()
#define ll long long
#define ld long double

void dbg_out(string s) { cerr << endl; }
template<typename H, typename... T>
void dbg_out(string s, H h, T... t){
    do{ cerr << s[0]; s = s.substr(1);
    } while (sz(s) && s[0] != ',');
    cerr << " = " << h;
    dbg_out(s, t...);
}

#ifdef DEBUG
#define dbg(...) dbg_out(#__VA_ARGS__, __VA_ARGS__)
#else
#define dbg(...) 42
#endif

void solution(){
    int n;
    cin >> n;
    for (int i=0; i<n; i++){
        
    }
}

signed main(){
    fastio;
    int cases=1;
    //cin >> cases;
    for (int i=0; i<cases; i++){
        solution();
    }
}
\end{lstlisting}
\section{Estruturas de Dados}
\subsection{Dsu}
\begin{lstlisting}
6a9 class DSU{
673   public:
1a8     int n;
f92     vector<int> p, r, s;
     
ea1     DSU(int nodes): p (nodes), r (nodes, 1), s (nodes, 1){
b67         n = nodes;
9e4         for (int i=0; i<n; i++) p[i] = i;
dc3     }
    	  /*
ac5     ll find(ll x){
eab         return (p[x] == x? x : p[x] = find(p[x]));
07f     }
c4c     */
57b     int find(int x){
42e         int root=x;
39f         while (root != p[root]) root = p[root];
133         int aux;
016         while (x != p[x]){
615             aux = p[x];
88a             p[x] = root;
835             x = aux;
8db         }
ea5         return x;
eb6     }
     
2e7     bool unite(int x, int y){
f08         x = find(x), y = find(y);
e65         if (x == y) return false;
7ce         if (r[x] > r[y]) swap(x, y);
b6a         if (r[x] == r[y]) r[y]++;
d16         s[y] += s[x];
ae9         p[x] = y;
15a         n--;
8a6         return true;
c37     }
     
082 };


\end{lstlisting}
\subsection{Segtree}
\begin{lstlisting}
853 template <class T> class Segtree{
673 public:
1a8     int n;
53f     T neutral;
22c     vector<T> t;
8c2     Segtree(int size, T neutral): neutral(neutral){
ba7         n = 1;
321         while (size > n) n <<= 1;
7ce         t.resize(2*n);
62b     }
    
1cb     T op(T a, T b){
534         return a + b;
4ed     }
    
6ab     T query(int l, int r, int idx, int a, int b){
41a         if (r < a || b < l) return neutral;
877         if (l <= a && b <= r) return t[idx];
e26         int m = (a+b)/2;
0f2         return op(query(l, r, 2*idx, a, m), query(l, r, 2*idx+1, m+1, b));
fd6     }
190     T query(int l, int r){ return query(l, r, 1, 0, n-1); }
    
671     void add(int idx, T v){
aab         if (idx < 0 || idx >= n) return;
cae         t[idx += n] += v; // aqui era pra ser op() invés de +=
4ef         for (idx >>= 1; idx; idx >>= 1) t[idx] = op(t[2*idx], t[2*idx+1]);
7ef     }
fdf };


\end{lstlisting}
\section{Grafos}
\subsection{Toposort}
Retorna as componentes na ordenação topológica reversa. O número de componentes é \textit{ncomp}.
Talvez o certo seja \textit{min(low[u], num[v])}, mas vou usar assim até achar um erro.

Estado em \textit{vis[]}:
\begin{itemize}
    \item \textbf{0}: não visitado
    \item \textbf{1}: finalizado
    \item \textbf{2}: ativo
\end{itemize}
\vspace{0.2cm}
\begin{lstlisting}
e3b vector<int> g [mx], comp (mx), low (mx), num (mx), vis (mx);
82b stack<int> vert;
22d int ti=0;
c02 int ncomp=0;
e0e void tarzan(int u){
873     low[u] = num[u] = ++ti;
10a     vis[u] = 2; // active
7a5     vert.push(u);
     
7b9     for (auto v: g[u]){
d90         if (!vis[v]){
5c1             tarzan(v);
ab6             low[u] = min(low[u], low[v]);
555         }
2d3         if (vis[v] == 2) low[u] = min(low[u], low[v]); // talvez o certo seja min(low[u], num[v]) mas vou usar assim até achar um erro
dd6     }
     
382     if (num[u] == low[u]){
a8f         ncomp++;
690         while (!vert.empty()){
08a             int v = vert.top(); vert.pop();
d6f             comp[v] = ncomp;
cca             vis[v] = 1;
163             if (u == v) break;
287         }
734     }
a7f }
\end{lstlisting}
\section{Teoria e Apêndices}
\subsection{Permutations}
\begin{itemize}
    \item O código para calcular a paridade da permutação está no caderno.
    \item \textbf{Face Value:} Use as definições de permutação de forma literal. Uma permutação de tamanho $N$ consiste em inteiros estritamente distintos de $1$ a $N$.
    \item \textbf{Divisibilidade:} Devido à limitação de limite, existem no máximo $\lfloor N/k \rfloor$ elementos na permutação que são múltiplos de $k$.
    \item \textbf{Validação Rápida:} Como o multiset de uma permutação válida de tamanho $M$ é sempre único $\{1, 2, \dots, M\}$, uma abordagem direta para checar se um subarray forma uma permutação é comparar os multisets utilizando algoritmos de \textbf{Hashing}.
\end{itemize}

\subsection{Lis}
\begin{itemize}
    \item Encontrar a LIS de um array reverso é equivalente a encontrar a \textbf{LDS} (Longest Decreasing Subsequence).
    \item Necessita de adaptação para \textit{longest non-increasing subsequence}. Para restaurar a subsequência, use um array auxiliar (vetor de pais).
    \item \textbf{Patience sorting:} Pensar no problema como formar pilhas, colocando o elemento atual na pilha com o menor topo que seja maior que ele. (Útil para modelagem geométrica ou de restrições).
    \item O número de LIS distintas pode ser calculado em $\mathcal{O}(N \log N)$ com SegTree.
    \item \textbf{LCS em Permutações:} A primeira permutação define o índice alvo para cada valor. Você substitui os valores da segunda permutação pelos índices correspondentes da primeira. A tarefa se reduz a achar a LIS do novo array, já que cada elemento aparece exatamente uma vez.
\end{itemize}

\end{document}